%%%%%%%%%%%%%%%%%%%%%%%%%%%%%%%%%%%%%%%%%%%%%%%%%%%%%%%%%%%
\begin{frame}[fragile]\frametitle{}
\begin{center}
{\Large Introduction to Deep NLP}
\end{center}
\end{frame}

%%%%%%%%%%%%%%%%%%%%%%%%%%%%%%%%%%%%%%%%%%%%%%%%%%%%%%%%%%%
\begin{frame}[fragile]\frametitle{Turing Test}

A human judge engages in a natural language conversation with two other parties, one a human and the other a machine; if the judge cannot reliably tell which is which, then the machine is said to pass the test.

\begin{center}
\includegraphics[width=0.4\linewidth,keepaspectratio]{turing}
\end{center}

Use of Deep Learning in Natural Language Processing is helping beat the Turing test, a mark of intelligence, ie Artificial Intelligence.
\end{frame}

%%%%%%%%%%%%%%%%%%%%%%%%%%%%%%%%%%%%%%%%%%%%%%%%%%%%%%%%%%%
\begin{frame}[fragile]\frametitle{NLP-AI}
	\begin{itemize}
	\item Natural Language Processing (Understanding-Generation) is an AI-Complete problem
	\item Most difficult AI problems are informally known as AI-complete or AI-hard
	\item To call a problem AI-complete reflects an attitude that it would not be solved by a simple specific algorithm.
	\end{itemize}

\tiny{(Ref: https://en.wikipedia.org/wiki/AI-complete)}
\end{frame}

%%%%%%%%%%%%%%%%%%%%%%%%%%%%%%%%%%%%%%%%%%%%%%%%%%%%%%%%%%%
\begin{frame}[fragile]\frametitle{NLP Tasks}
	\begin{itemize}
	\item Part Of Speech Tagging: Assign part-of-speech to each word.
	\item Named Entity Recognition: Recognize people, places, etc. in a sentence.
	\item Language Modeling: Predict next word, Generate natural sentences.
	\item Translation: Translate a sentence into another language.
	\item Summarization: Summarize a paragraph in new words.
	\end{itemize}

\tiny{(Ref: How Deep Learning Quietly Revolutionized NLP - Lukasz Kaiser, Google Brain)}
\end{frame}

%%%%%%%%%%%%%%%%%%%%%%%%%%%%%%%%%%%%%%%%%%%%%%%%%%%%%%%%%%%
\begin{frame}[fragile]\frametitle{What is Deep NLP}
\begin{center}
\includegraphics[width=0.8\linewidth,keepaspectratio]{nlp4}

\tiny{(Ref: Deep Learning and NLP A-Z - Kirill Eremenko)}

\tiny{(Note: Size is not indicative of importance)}
\end{center}

	\begin{itemize}
	\item Green part is NLP (rule based, linguistic)
	\item Blue part is Deep Learning not applied to NLP
	\item Purple is Deep NLP (DNLP), NN applied for NLP use cases
	\item Seq2Seq is a heavily used technique of DNLP for sequence to sequence modeling, eg Translation, Q \& A, etc.
	\end{itemize}
\end{frame}

%%%%%%%%%%%%%%%%%%%%%%%%%%%%%%%%%%%%%%%%%%%%%%%%%%%%%%%%%%%
\begin{frame}[fragile]\frametitle{Machine Learning: Hand-Crafted Features}
\begin{center}
\includegraphics[width=0.8\linewidth,keepaspectratio]{dnlp1}

\tiny{(Ref: Deep Learning for Natural Language Processing - Sihem Romdhani)}
\end{center}

Hand crafted features are needed in Machine Learning. E.g. for Spam Detection, features could be presence of BIG \$ amounts, FROM country, etc.
\end{frame}

%%%%%%%%%%%%%%%%%%%%%%%%%%%%%%%%%%%%%%%%%%%%%%%%%%%%%%%%%%%
\begin{frame}[fragile]\frametitle{Deep Learning: Learned Features}
\begin{center}
\includegraphics[width=0.6\linewidth,keepaspectratio]{dnlp2}

\tiny{(Ref: Deep Learning for Natural Language Processing - Sihem Romdhani)}
\end{center}

Hand crafted features are NOT needed in Deep Learning. E.g. for Object Detection, CNNs come up with their own features like edges, parts, etc.
\end{frame}

%%%%%%%%%%%%%%%%%%%%%%%%%%%%%%%%%%%%%%%%%%%%%%%%%%%%%%%%%%%
\begin{frame}[fragile]\frametitle{Need Numbers, Not Words!}
How can we represent words?
\begin{center}
\includegraphics[width=0.8\linewidth,keepaspectratio]{dnlp3}

\tiny{(Ref: Deep Learning for Natural Language Processing - Sihem Romdhani)}
\end{center}

Words, sentences, paragraphs, pages, documents all need to be represented by numbers, to be able to be fed to an ML/DL algorithm.
\end{frame}

%%%%%%%%%%%%%%%%%%%%%%%%%%%%%%%%%%%%%%%%%%%%%%%%%%%%%%%%%%%
\begin{frame}[fragile]\frametitle{Traditional Encoding and Its Challenges}
With Vocabulary as the full set representing the whole corpus:
	\begin{itemize}
	\item Bag-of-Words: Count of occurrences of a word in a document
	\item TF-IDF: Measures importance of a word in a document relative to the corpus
	\item[] \textbf{Challenges:} Sparse inputs, and context is lost in encoding.
	\end{itemize}

\begin{center}
\includegraphics[width=0.7\linewidth,keepaspectratio]{dnlp9}

\tiny{(Ref: Deep Learning for Natural Language Processing - Bargava, Amit Kapoor)}
\end{center}
\end{frame}

%%%%%%%%%%%%%%%%%%%%%%%%%%%%%%%%%%%%%%%%%%%%%%%%%%%%%%%%%%%
\begin{frame}[fragile]\frametitle{Embeddings}
It can not be just any set of numbers representing words. Similar words should have similar numbers -- context defines similarity, so we infer the meaning of a word from the company it keeps.

\begin{center}
\includegraphics[width=0.5\linewidth,keepaspectratio]{dnlp4}

\tiny{(Ref: Deep Learning for Natural Language Processing - Sihem Romdhani)}
\end{center}
\end{frame}

%%%%%%%%%%%%%%%%%%%%%%%%%%%%%%%%%%%%%%%%%%%%%%%%%%%%%%%%%%%
\begin{frame}[fragile]\frametitle{Word to Vector}
Similar words end up nearby in vector space; dissimilar words are further away.

\begin{center}
\includegraphics[width=0.6\linewidth,keepaspectratio]{dnlp7}

\tiny{(Ref: Deep Learning for Natural Language Processing - Sihem Romdhani)}
\end{center}
\end{frame}

%%%%%%%%%%%%%%%%%%%%%%%%%%%%%%%%%%%%%%%%%%%%%%%%%%%%%%%%%%%
\begin{frame}[fragile]\frametitle{Vector Algebra Possible}

\begin{center}
\includegraphics[width=\linewidth,keepaspectratio]{dnlp10}

\tiny{(Ref: https://www.tensorflow.org/versions/r0.12/tutorials/word2vec/index.html)}
\end{center}
\end{frame}

%%%%%%%%%%%%%%%%%%%%%%%%%%%%%%%%%%%%%%%%%%%%%%%%%%%%%%%%%%%
\begin{frame}[fragile]\frametitle{Reasons for Applying Deep Learning to NLP}
	\begin{itemize}
	\item \textbf{Automatic representation learning} -- no hand-crafted features needed.
	\item \textbf{Learning from unlabeled data} -- methods like Skip-Gram and CBOW generate word2vec by making labels out of unlabeled data, unlike traditional ML-based NLP which needs labeled training data.
	\item \textbf{Sequential and contextual} -- human language is inherently sequential; RNNs serve that purpose well.
	\end{itemize}

\begin{center}
\includegraphics[width=0.7\linewidth,keepaspectratio]{dnlp11}

\tiny{(Ref: A not-so-short introduction to Deep Learning NLP - Francesco Gadaleta, PhD)}
\end{center}
\end{frame}

% %%%%%%%%%%%%%%%%%%%%%%%%%%%%%%%%%%%%%%%%%%%%%%%%%%%%%%%%%%%
% \begin{frame}[fragile]\frametitle{Seq2Seq Architecture}
%
% \begin{center}
% \includegraphics[width=0.8\linewidth,keepaspectratio]{nlp14}
%
% \tiny{(Ref: Deep Learning and NLP A-Z - Kirill Eremenko)}
% \end{center}
%
% During training, the Encoder is fed Questions and the Decoder is fed Answers; weights in the gates and hidden states get settled. During testing, for each input sequence the encoder produces a combined vector, and the decoder takes this and starts spitting out words one by one, probabilistically.
% \end{frame}

% %%%%%%%%%%%%%%%%%%%%%%%%%%%%%%%%%%%%%%%%%%%%%%%%%%%%%%%%%%%
% \begin{frame}[fragile]\frametitle{Attention Mechanism}
% 	\begin{itemize}
% 	\item A single meaning/context vector ($h_n$), the final encoder vector, to carry a whole sentence is too much to expect.
% 	\item We need the luxury of looking back at the source, whenever needed.
% 	\item Through training, the decoder learns, for a particular output word, which of the input words need to be looked at -- that source word gets the highest weight.
% 	\end{itemize}
%
% \begin{center}
% \includegraphics[width=0.6\linewidth,keepaspectratio]{nlp20}
%
% \tiny{(Ref: Deep Learning and NLP A-Z - Kirill Eremenko)}
% \end{center}
% \end{frame}

% %%%%%%%%%%%%%%%%%%%%%%%%%%%%%%%%%%%%%%%%%%%%%%%%%%%%%%%%%%%
% \begin{frame}[fragile]\frametitle{Decoding: Greedy vs Beam Search}
% \textbf{Greedy Decoding}: at each time step, pick the word with max probability out of the whole vocabulary, feed it back in, and repeat until EOS is produced.
%
% \textbf{Beam Search Decoding}: instead of only the top-1 word, keep the top-$k$ (eg top 3) candidate sequences at each step, like a tree, and pick the winner by max joint probability.
%
% \begin{center}
% \includegraphics[width=0.6\linewidth,keepaspectratio]{nlp18}
%
% \tiny{(Ref: Deep Learning and NLP A-Z - Kirill Eremenko)}
% \end{center}
% \end{frame}

% %%%%%%%%%%%%%%%%%%%%%%%%%%%%%%%%%%%%%%%%%%%%%%%%%%%%%%%%%%%
% \begin{frame}[fragile]\frametitle{Named Entity Recognition: Traditional vs Deep}
% \textbf{Traditional}: Conditional Random Fields (CRFs), a ``classic'' machine learning approach.
%
% \textbf{Deep}: LSTM / Bi-LSTM sequence labelling.
%
% \begin{center}
% \includegraphics[width=0.6\linewidth,keepaspectratio]{dnlp19}
%
% \tiny{(Ref: Deep Learning for NLU - Dr. David Talby)}
% \end{center}
% \end{frame}

%%%%%%%%%%%%%%%%%%%%%%%%%%%%%%%%%%%%%%%%%%%%%%%%%%%%%%%%%%%
\begin{frame}[fragile]\frametitle{Conclusion}
	\begin{itemize}
	\item Deep learning has simplified feature engineering in many cases (it certainly hasn't removed it).
	\item Less feature engineering is leading to more complex machine learning architectures.
	\item Most of the time, these model architectures are as specific to a given task as feature engineering used to be.
	\item The job of the data scientist will stay sexy for a while (keep your fingers crossed on this one).
	\end{itemize}
\end{frame}
